\section*{Chapter \ref{ch:s1:share-classes-holding-companies-closed-end-funds-spacs} --- Share Classes, Holding Companies, Closed-End Funds and SPACs}

\begin{solution}{exo:s1:share-classes-holding-companies-closed-end-funds-spacs:1}
$\ln 0.85 = -16.3\%$, $\ln 0.90 = -10.5\%$, $\ln 0.95 = -5.1\%$.
\end{solution}

\begin{solution}{exo:s1:share-classes-holding-companies-closed-end-funds-spacs:2}
$\ln 0.95 - \ln 0.85 - 0.013 = 11.1\% - 1.3\% = 9.8\%$.
\end{solution}

\begin{solution}{exo:s1:share-classes-holding-companies-closed-end-funds-spacs:3}
$10 \times e^{0.04} = \$10.41$; the floor return is $10.41/9.80 - 1 = 6.2\%$.
\end{solution}

\begin{solution}{exo:s1:share-classes-holding-companies-closed-end-funds-spacs:4}
$0.5^{1/2} = 70.7\%$ of the gap remains after a year. The expected narrowing is to the average discount, not to zero, and slowly; carry runs the whole time,
so the expected gain per year is small and the time to reach the exit level long.
\end{solution}

\begin{solution}{exo:s1:share-classes-holding-companies-closed-end-funds-spacs:5}
Share classes are claims on the same company and differ only in rights and liquidity, so their spread is small and reverts fast. A fund's shares cannot be
redeemed, its portfolio may be hard to replicate, its expenses run while it is held, and its discount moves with investor sentiment across funds: nothing
forces it closed.
\end{solution}

\begin{solution}{exo:s1:share-classes-holding-companies-closed-end-funds-spacs:6}
At the vote the holder can redeem for the \omterm{def:s1:share-classes-holding-companies-closed-end-funds-spacs:spac}{trust value} or keep the merged share: the payoff is the maximum of the two, a floor plus a call on the merged
company. The floor fails if the trust does not hold what it should (it earns less, or the SPAC's terms change through extensions), or if the shares cannot
be redeemed as expected.
\end{solution}

\begin{solution}{exo:s1:share-classes-holding-companies-closed-end-funds-spacs:7}
At 0.1 catalysts a year the trade earns 2.23\% a year and closes in a median 2.13 years; without a catalyst 0.87\% and 2.96 years; at 0.3, 5.04\% and 1.37
years. Going from none to 0.3 adds about 4.2 points a year on the capital committed: that is the most a catalyst is worth to the trade, before the costs of
producing it.
\end{solution}

\begin{solution}{exo:s1:share-classes-holding-companies-closed-end-funds-spacs:8}
Nothing makes the fund's price rise to its asset value: the hedged trade earns only the change in the discount, which can widen, less the fund's expenses and
the hedge's costs. The expected narrowing is to the fund's average discount; without a catalyst the model trade earns 1.9\% over a median three years and
loses in 40\% of the paths.
\end{solution}

\begin{solution}{pb:s1:share-classes-holding-companies-closed-end-funds-spacs:1}
\begin{enumerate}
\item Two listed classes with the same cash-flow claim and different rights; the shortfall of a holding company's value from its stakes; the shortfall of a
closed-end fund's price from its net asset value.
\item Voting and liquidity differences, taxes and costs, the inability to redeem; nothing forces them closed cheaply.
\item Share classes are real twins: a small stationary spread, which chapter~4's \omterm{def:s1:pairs-and-baskets:distance}{distance method} traded at a Sharpe ratio of 5.
\item Fund discounts move together with sentiment and narrow when small stocks do well; deviations are larger where arbitrage is costlier.
\item A log discount from $-16.3\%$ reverting toward $-10.5\%$ with a two-year half-life and 8\% volatility; long the fund, short its portfolio, 1.3\% carry,
exit at a 5\% discount, a catalyst or three years.
\item 1.9\% per trade, 0.87\% a year, median three years, half closed, 40\% losing.
\item 2.23\% and 5.04\% a year; median 2.13 and 1.37 years; 63\% and 79\% closed within three years.
\item Without a catalyst and with slow reversion, a wide discount is a poor trade: 0.94\% a year, 2\% closed within three years.
\item A listed shell with a trust, a deadline and redemption rights; the redemption amount per share.
\item \$10.41 at the vote, a 6.2\% return from \$9.80.
\item Net cash per share at merger far below \$10 (median \$5.70 in 2019--2020) and steep post-merger losses for holders.
\item Redeeming below the trust (74\% of paths) gives 20.2\% on average; holding always loses 13.9\% on average and 27.8\% at the median.
\item \textbf{Named result.} Bought at 85 cents with no catalyst, the discount trade earns 0.87\% a year and closes in a median 2.96 years (half the paths not
closing within three); with catalysts at 0.3 a year, 5.04\% a year and a median 1.37 years.
\item A price tied by a structure to an observable value, a gap no one can close cheaply, carry while waiting.
\item Because the activist owns the catalyst.
\item Enhanced disclosure of conflicts, sponsor compensation, dilution and the target, and rules on projections.
\item With a proxy (an index or a replicating basket) and accept the tracking error, sized accordingly.
\item The holding-company or fund discount with a catalyst in view.
\item Discounts widening together when sentiment turns.
\item A gap between a price and the value it is tied to, and the time it takes to close.
\end{enumerate}
\end{solution}

\begin{solution}{iq:s1:share-classes-holding-companies-closed-end-funds-spacs:1}
Its shares cannot be redeemed at net asset value, so nothing forces the price to it; expenses, taxes, illiquid holdings and investor sentiment keep a
discount, and arbitrage against it is costly and slow.
\end{solution}

\begin{solution}{iq:s1:share-classes-holding-companies-closed-end-funds-spacs:2}
Long the cheap class, short the rich one, when the spread is wide relative to its history and to the value of the voting difference; small positions, many
trades, and closing before events that can change the value of the vote.
\end{solution}

\begin{solution}{iq:s1:share-classes-holding-companies-closed-end-funds-spacs:3}
As the discounted \omterm{def:s1:share-classes-holding-companies-closed-end-funds-spacs:spac}{trust value} plus a call on the merged company's share struck at the \omterm{def:s1:share-classes-holding-companies-closed-end-funds-spacs:spac}{trust value}, with warrants valued separately; the call's value depends
on the likely merger, its dilution and the post-merger price distribution.
\end{solution}

\begin{solution}{iq:s1:share-classes-holding-companies-closed-end-funds-spacs:4}
Discounts widening together (common sentiment), hedges that do not track the portfolios, expenses and carry over long holding periods, and illiquidity in
the funds themselves.
\end{solution}

\begin{solution}{iq:s1:share-classes-holding-companies-closed-end-funds-spacs:5}
A catalyst: a planned sale or distribution of stakes, a change of control or governance, an activist with influence, or a tax change; and a hedge for the
listed stakes, so that the position is a bet on the discount alone.
\end{solution}

\begin{solution}{iq:s1:share-classes-holding-companies-closed-end-funds-spacs:6}
Before the catalyst, $E[d_t] = \bar d + (d_0 - \bar d)e^{-\kappa t}$; with probability $e^{-\lambda t}$ no catalyst has arrived. So $E[d_t] = e^{-\lambda t}(\bar d +
(d_0 - \bar d)e^{-\kappa t}) + (1 - e^{-\lambda t})d_c$, and the expected gain of the trade held to $t$ is $E[d_t] - d_0$ less carry.
\end{solution}
